\documentclass[11pt,a4paper]{article}
\usepackage[margin=20mm,top=19mm,bottom=20mm]{geometry}
\usepackage{fontspec,amsmath,amssymb,xcolor,enumitem,fancyhdr}
\setmainfont{NotoSansKR-Regular.ttf}[Path=../../../common/assets/fonts/,BoldFont=NotoSansKR-Semibold.ttf]
\setmonofont{Latin Modern Mono}
\XeTeXlinebreaklocale "ko"
\XeTeXlinebreakskip=0pt plus 1pt
\definecolor{ink}{HTML}{203B50}
\definecolor{rulegray}{HTML}{D2D9DF}
\pagestyle{fancy}\fancyhf{}
\fancyhead[L]{\small 머신러닝 개념 확인}\fancyhead[R]{\small W2}
\fancyfoot[C]{\small\thepage\ / 4}
\renewcommand{\headrulewidth}{0.3pt}
\setlength{\headheight}{15pt}
\setlength{\parindent}{0pt}\setlength{\parskip}{5pt}
\linespread{1.2}
\newcommand{\parttitle}[2]{{\large\bfseries\color{ink} #1}\par{\small #2}\par\vspace{5mm}}
\newcommand{\question}[2]{\textbf{#1.} #2\par}
\newcommand{\answerlines}[1]{\foreachline{#1}}
\newcount\linecount
\newcommand{\foreachline}[1]{\linecount=0\loop\ifnum\linecount<#1\vspace{8mm}\par{\color{rulegray}\hrule}\advance\linecount by1\repeat}
\begin{document}

{\LARGE\bfseries\color{ink} W2 · 여러 특성을 이용한 회귀}\par
{\small 다중 선형 회귀 · 벡터화 · 특성 스케일링 · 다항식 특성}\par
\vspace{3mm}이름\enspace\rule{43mm}{0.3pt}\hfill 날짜\enspace\rule{35mm}{0.3pt}\par\vspace{3mm}
\parttitle{1. 개념 확인}{각 문항에서 가장 적절한 답 하나를 고르세요.}
\question{1}{집값을 예측할 때 넓이뿐 아니라 방 개수와 집의 나이도 사용하려 합니다. 가장 적절한 설명은 무엇일까요?}
\begin{enumerate}[label=(\arabic*),leftmargin=9mm,itemsep=0pt,topsep=2pt,parsep=0pt]
\item 정답 가격을 세 번 복사한다
\item 세 특성 중 반드시 하나만 남긴다
\item 비용 함수를 없앤다
\item 집 하나를 여러 입력 특성으로 표현한다
\end{enumerate}\vfill
\question{2}{$\hat y=w_1x_1+w_2x_2+b$에서 각 가중치의 역할은 무엇일까요?}
\begin{enumerate}[label=(\arabic*),leftmargin=9mm,itemsep=0pt,topsep=2pt,parsep=0pt]
\item 해당 특성이 예측값에 기여하는 정도를 조절한다
\item 각 특성의 이름을 정한다
\item 데이터의 행 개수만 정한다
\item 실제 정답을 임의로 바꾼다
\end{enumerate}\vfill
\question{3}{$x=(2,3)$, $w=(4,1)$, $b=1$이면 $\hat y=w^Tx+b$는 얼마일까요?}
\begin{enumerate}[label=(\arabic*),leftmargin=9mm,itemsep=0pt,topsep=2pt,parsep=0pt]
\item $10$
\item $11$
\item $12$
\item $14$
\end{enumerate}\vfill
\question{4}{벡터화를 사용하는 주된 이유는 무엇일까요?}
\begin{enumerate}[label=(\arabic*),leftmargin=9mm,itemsep=0pt,topsep=2pt,parsep=0pt]
\item 다른 정답으로 학습하기 위해
\item 같은 계산을 배열 연산으로 간결하고 효율적으로 처리하기 위해
\item 모든 특성을 하나의 상수로 바꾸기 위해
\item 모델의 오차를 무조건 0으로 만들기 위해
\end{enumerate}\vfill
\question{5}{넓이는 수천 단위이고 방 개수는 한 자릿수입니다. 경사 하강법 전에 특성 스케일링을 하는 이유로 가장 적절한 것은 무엇일까요?}
\begin{enumerate}[label=(\arabic*),leftmargin=9mm,itemsep=0pt,topsep=2pt,parsep=0pt]
\item 실제 집값을 모두 같게 만들기 위해
\item 작은 값을 가진 특성을 삭제하기 위해
\item 넓이와 방 개수를 같은 의미로 바꾸기 위해
\item 특성별 수치 범위 차이를 줄여 학습이 원활하도록 돕기 위해
\end{enumerate}\vfill
\newpage
\parttitle{1. 개념 확인 · 계속}{각 문항에서 가장 적절한 답 하나를 고르세요.}
\question{6}{평균이 $10$, 표준편차가 $2$인 특성에서 값 $14$를 $(x-\mu)/\sigma$로 표준화하면 얼마일까요?}
\begin{enumerate}[label=(\arabic*),leftmargin=9mm,itemsep=0pt,topsep=2pt,parsep=0pt]
\item $2$
\item $4$
\item $7$
\item $12$
\end{enumerate}\vfill
\question{7}{새 집의 특성을 표준화해 예측하려 합니다. 어떤 평균과 표준편차를 사용해야 할까요?}
\begin{enumerate}[label=(\arabic*),leftmargin=9mm,itemsep=0pt,topsep=2pt,parsep=0pt]
\item 새 집 한 채만으로 구한 값
\item 예측할 때마다 무작위로 정한 값
\item 학습 데이터에서 구해 저장한 값
\item 정답 가격의 평균과 표준편차
\end{enumerate}\vfill
\question{8}{학습 중 비용이 계속 커지고 크게 흔들립니다. 먼저 시도할 만한 조치는 무엇일까요?}
\begin{enumerate}[label=(\arabic*),leftmargin=9mm,itemsep=0pt,topsep=2pt,parsep=0pt]
\item 학습률을 더 크게 올린다
\item 학습률을 낮추고 비용 변화를 확인한다
\item 실제 정답을 전부 0으로 바꾼다
\item 비용을 기록하지 않는다
\end{enumerate}\vfill
\question{9}{넓이 $x$에 따라 가격이 곡선 형태로 변합니다. 이번 강의의 방법으로 시도할 수 있는 것은 무엇일까요?}
\begin{enumerate}[label=(\arabic*),leftmargin=9mm,itemsep=0pt,topsep=2pt,parsep=0pt]
\item 모든 입력을 1로 바꾼다
\item 모든 가중치를 0으로 고정한다
\item 같은 입력의 정답을 무작위로 바꾼다
\item $x^2$ 같은 새 특성을 추가한다
\end{enumerate}\vfill
\question{10}{$\hat y=w_1x+w_2x^2+b$가 입력 $x$에 대해서는 곡선이어도 선형 회귀로 다룰 수 있는 이유는 무엇일까요?}
\begin{enumerate}[label=(\arabic*),leftmargin=9mm,itemsep=0pt,topsep=2pt,parsep=0pt]
\item 매개변수 $w_1,w_2,b$가 선형적으로 결합되기 때문이다
\item 모든 곡선이 직선이기 때문이다
\item 입력 값이 반드시 0과 1이기 때문이다
\item 가중치를 학습하지 않기 때문이다
\end{enumerate}\vfill
\newpage
\parttitle{2. 머신러닝 사이언티스트의 점검}{11번 · 의사코드를 읽고, 모델이 의도대로 작동하지 않는 원인을 찾으세요.}
넓이와 방 개수를 이용해 집값을 예측합니다. 학습할 때는 각 특성을 표준화했고, 학습 오차도 작았습니다. 그런데 새 집을 예측하면 터무니없는 가격이 나옵니다. 아래에서 mu와 sigma는 특성별 평균과 표준편차이며, sigma의 값은 모두 0보다 큽니다. 학습 함수는 올바르게 작동합니다.
\par\vspace{3mm}
{\small\ttfamily\begin{tabular}{r l}
1 & \hspace*{0.0mm}mu = column\_mean(X\_train) \\
2 & \hspace*{0.0mm}sigma = column\_std(X\_train) \\
3 & \hspace*{0.0mm}X\_scaled = (X\_train - mu) / sigma \\
4 & \hspace*{0.0mm}w, b = fit\_linear(X\_scaled, y\_train) \\
5 & \hspace*{0.0mm}x\_new = [1500, 3] \\
6 & \hspace*{0.0mm}prediction = dot(x\_new, w) + b \\
\end{tabular}}\par\vspace{5mm}
\textbf{11(1)} 학습할 때의 입력과 예측할 때의 입력을 비교해, 빠진 처리가 무엇인지 설명하세요.
\answerlines{5}
\vspace{6mm}
\textbf{11(2)} 필요한 처리를 한 줄 추가하고 마지막 줄을 한 줄 수정하세요. 어느 데이터에서 구한 통계값을 써야 하는지도 적으세요.
\answerlines{5}
\vspace{6mm}
\newpage
\parttitle{3. 열린 질문}{12번 · 정해진 정답은 없습니다. 제안한 방법과 그 이유가 드러나면 충분합니다.}
문장을 같은 길이의 숫자 목록으로 바꾸는 도구가 있습니다. 여기에 문장 길이와 느낌표 개수도 특성으로 추가할 수 있습니다. 비슷한 문의 내용을 찾는 시스템을 만든다면, 어떤 특성을 사용하고 어떻게 비교하겠습니까?
\par\vspace{3mm}
특성들의 숫자 범위가 많이 다르면 어떻게 할지, 특성을 추가하기 전과 후에 무엇을 비교해 볼지 설명하세요. 도구가 만든 숫자가 어떤 의미를 갖는지는 직접 확인해야 한다고 생각해도 좋습니다.
\par\vspace{3mm}
{\small 수식이나 코드를 반드시 쓸 필요는 없습니다. 글, 그림, 간단한 순서도로 답해도 좋습니다.}
\answerlines{15}
\vfill{\footnotesize 이 문제지는 Machine Learning Specialization Course 1의 해당 주차 개념을 바탕으로 새로 작성한 학습용 자료입니다.}
\end{document}